\documentclass[12pt]{article}

\usepackage{setspace}
\onehalfspace

%\usepackage[utf8]{inputenc}
%\usepackage{t1enc}
%\usepackage[T1]{fontenc}

\usepackage{amsmath}

\usepackage{moodle}

\DeclareMathOperator{\Exp}{Exp}
\DeclareMathOperator{\Norm}{{\cal N}}
\DeclareMathOperator{\Binom}{Binom}
\DeclareMathOperator{\E}{E}
\DeclareMathOperator{\D}{D}
\renewcommand{\P}{\operatorname{P}}
%\DeclareMathOperator{\cov}{cov}
\DeclareMathOperator{\corr}{corr}

\DeclareMathOperator{\F}{F} 
\renewcommand{\d}{\operatorname{d\! }} %integral \! a nyero


\DeclareMathOperator{\R}{\mathbb{R}}
%\DeclareMathOperator{\Norm}{{\cal N}}
\DeclareMathOperator{\Unif}{{\cal U}}
\DeclareMathOperator{\Poi}{{Poisson}}


\DeclareMathOperator{\cov}{cov}





\begin{document}

\begin{quiz}{gyakorló}

\begin{multi}{felt9}
Ha $A_1,\ldots,A_n$ teljes eseményrendszer pozitív valségű tagokkal és $\P(B)>0$ akkor minden $m$-re:
$$
\P(A_m|B)=\frac{\P(B|A_m)\P(A_m)}{\sum_{k=1}^n \P(B|A_k)\P(A_k)}
$$
\item* igaz
\item hamis
\end{multi}
\begin{multi}{dvv4}
Egy $\xi$ valségi változó az $x_1,\ldots x_n$ értékeket veheti fel. Ekkor a várható értéke:
$$
\E(\xi)=\frac{\sum_{k=1}^n x_k}{n}
$$
\item igaz
\item* hamis
\end{multi}
\begin{multi}{eofv2}
Azt mondjuk, hogy egy $F:\R\to \R$ eloszlásfüggvény, ha $F\ge 0$ és 
$$
\int_{-\infty}^{+\infty}F(x)\d x=1
$$
\item igaz
\item* hamis
\end{multi}
\begin{multi}{impdvv5}
Azt mondjuk, hogy $\xi\sim \Binom(n,p)$, ha 
$$
\P(\xi=k)=\binom{n}{k}p^k(1-p)^{n-k}\ \ \ k=0\ldots n
$$
valamilyen $n>0$ és $p\in[0,1]$ esetén.
\item* igaz
\item hamis
\end{multi}
\begin{multi}{impdvv7}
Azt mondjuk, hogy $\xi\sim \Poi(\lambda)$, ha 
$$
\P(\xi=k)=e^{-\lambda}\frac{\lambda^k}{k!}\ \ \ (k=0,1,2\ldots )
$$ 
valamilyen $\lambda>0$ esetén.
\item* igaz
\item hamis
\end{multi}
\begin{multi}{eofv9}
Bármely $\xi$-re és $a$-ra $F_{\xi}(a)+\P(\xi\ge a)=1$.
\item* igaz
\item hamis
\end{multi}
\begin{multi}{dvv14}
$\xi$ és $\eta$ független, véges értékkészletű valségi változókra:
$$
\D^2(\xi+\eta)=\D^2(\xi)+\D^2(\eta)
$$
\item* igaz
\item hamis
\end{multi}
\begin{multi}{esem4}
Az $\overline{A} + \overline{B}$ maga után vonja az $\overline{A\cdot B}$ eseményt.
\item* igaz
\item hamis
\end{multi}
\begin{multi}{dvv3}
Egy $\xi$ valségi változó az $x_1,\ldots x_n$ értékeket veheti fel. Ekkor a várható értéke:
$$
\E(\xi)=\sum_{k=1}^n \P(\xi=x_k)x_k
$$
\item* igaz
\item hamis
\end{multi}
\begin{multi}{flenvv6}
Ha $\xi$ és $\eta$ valségi változóknak létezik a szórása, akkor a 
kovarianciájuk
$$
\cov(\xi,\eta)=\E(\xi\eta)-\E(\xi)\E(\eta)
$$
\item* igaz
\item hamis
\end{multi}
\begin{multi}{impfvv10}
Ha $\xi\sim \Exp(\lambda)$ akkor a sűrűségfüggvénye 
$$
f_{\xi}(x)=
\begin{cases}
1-e^{-\lambda x}& \ \ x>0\\
0 & \text{máskor} 
\end{cases}
$$
\item igaz
\item* hamis
\end{multi}
\begin{multi}{impfvv9}
Ha $\xi\sim \Exp(\lambda)$ akkor a sűrűségfüggvénye 
$$
f_{\xi}(x)=
\begin{cases}
\lambda e^{-\lambda x}& \ \ x>0\\
0 & \text{máskor} 
\end{cases}
$$
\item* igaz
\item hamis
\end{multi}
\begin{multi}{flenvv5}
Ha $\xi$ és $\eta$ diszkrét valségi változóknak létezik a szórása és nemnulla, akkor a 
korrelációjuk
$$
\corr(\xi,\eta)=\E(\xi\eta)-\E(\xi)\E(\eta)
$$
\item igaz
\item* hamis
\end{multi}
\begin{multi}{sfv9}
Ha $\xi$-nek $f$ a sűrűségfüggvénye és $F$ az eloszlásfüggvénye, akkor bármely $b$-re:
$$
F(b)=\int_{\infty}^{b}f(x)\d x
$$
\item* igaz
\item hamis
\end{multi}
\begin{multi}{val6}
Bármely esemény valsége kiszámolható a $\frac{\text{kedvező}}{\text{összes}}$ képlettel, ha az elemi események 
egyforma valségűek.
\item* igaz
\item hamis
\end{multi}
\begin{multi}{felt5}
$A_1,\ldots,A_n$ teljes eseményrendszer, ha $\sum A_k=\Omega$ és a tagok páronként kizáróak.
\item* igaz
\item hamis
\end{multi}
\begin{multi}{val7}
Bármely $C$-re $\P(C)=\sum_{\omega_k\in C} \P( \{\omega_k \})$.
\item* igaz
\item hamis
\end{multi}
\begin{multi}{felt6}
$A_1,\ldots,A_n$ teljes eseményrendszer, ha $\sum A_k=\Omega$ és a tagok páronként függetlenek.
\item igaz
\item* hamis
\end{multi}
\begin{multi}{dvv10}
Egy $\xi$ valségi változó az $x_1,\ldots x_n$ értékeket veheti fel. 
Ekkor a szórásnégyzete:
$$
\D^2(\xi)=\E(\xi^2)-\E(\xi)^2
$$
\item* igaz
\item hamis
\end{multi}
\begin{multi}{esem2}
Az $\overline{A+B}$ és $\overline{A} \cdot \overline{B}$ események kizárják egymást.
\item igaz
\item* hamis
\end{multi}
\begin{multi}{eofv7}
Bármely $\xi$-re és $a<b$-re $\P(a\le \xi <b)=F_{\xi}(b)-F_{\xi}(a)$.
\item* igaz
\item hamis
\end{multi}
\begin{multi}{impfvv5}
Ha $\xi\sim \Unif(0,1)$ akkor minden $a,b$-re $(b-a)\xi+a\sim \Unif(a,b)$.
\item* igaz
\item hamis
\end{multi}
\begin{multi}{sfv3}
Azt mondjuk, hogy egy $f:\R\to \R$ sűrűségfüggvény 
$$
\int_{-\infty}^{+\infty}f(x)\d x\ge 0
$$
\item igaz
\item* hamis
\end{multi}
\begin{multi}{sfv8}
Ha van $\xi$-nek sűrűségfüggvénye, akkor bármely $b$-re:
$$
\P(\xi<b)=f(b)
$$
\item igaz
\item* hamis
\end{multi}
\begin{multi}{esem5}
Ha $\P(A)=1$ akkor $A$ bekövetkezhet.
\item* igaz
\item hamis
\end{multi}
\end{quiz}

\end{document}

